\textbf{Scale.} One easy dataset (569 patients; LR reaches AUROC 0.993), 10 seeds of one 70/30 split protocol, no external validation and no tuning. Discrimination has little room to improve, so a null is expected; harder tasks may differ.

\textbf{Few positives.} At 1:20 and 1:50 the models are trained on 12 and 5 malignant cases, and we evaluate on all held-out patients with positives down-weighted to the target prevalence. Brier score, AUPRC, calibration slope and intercept are then dominated by the 107 negatives and are noisy; GB's $|a|$ of 4 to 6 logits may reflect a few saturated errors. The weighted evaluation is not equivalent to a genuine 1:50 cohort.

\textbf{Statistics.} Tests are unadjusted paired $t$-tests over 10 seeds that share one dataset, and \texttt{results/tables} holds several hundred of them; isolated $p<0.05$ cells are not confirmed effects. The signed-intercept, signed-slope and operating-point analyses were not registered and were added after seeing the data; we rely on them only where all LR cells agree. The registered H4 threshold turned out to be nearly vacuous, which we report rather than change.

\textbf{Implementation.} SMOTE, ROS and the offset are reimplemented and were checked only through the run logs (e.g.\ training-set sizes), not against imbalanced-learn. One SMOTE variant ($k=5$) and fixed GB hyperparameters were used; tuned or regularised boosting, undersampling, cost-sensitive learning and post-hoc recalibration (Platt, isotonic) are out of scope, so we cannot say whether a corrected and then recalibrated model would match no correction. The mechanisms offered above are untested. A full study would add several datasets, nested tuning, calibration curves and genuinely imbalanced cohorts.
